\begin{topic}[id=cra_embedded_iot,
             difficulty={Mittel},
             type={Technische Regulierungs- und Systemanalyse},
             prerequisites={Grundlagen der IT-Sicherheit, eingebettete Systeme, vernetzte Geraete},
             practice={gut moeglich},
             owner={Dozententeam},
             updated={2026-04-13},
             tags={ CRA, Embedded Security, IoT-Sicherheit, Security by Design, Schwachstellenmanagement, OTA-Updates, SBOM, Incident Reporting }]{Cyber Resilience Act fuer vernetzte Embedded- und IoT-Produkte}

\TopicDescription{Der Cyber Resilience Act (CRA) betrifft nicht nur klassische Software, sondern auch vernetzte Embedded- und IoT-Produkte, deren Sicherheitsniveau ueber den gesamten Produktlebenszyklus nachvollziehbar abgesichert werden muss. Fuer ein seminarfaehiges Thema steht daher nicht die juristische Detailauslegung im Vordergrund, sondern die technische Uebersetzung regulatorischer Vorgaben in Architektur-, Entwicklungs- und Wartungsprozesse. Untersucht werden Security-by-Design und Security-by-Default fuer Geraetefunktionen, Schnittstellen und Standardkonfigurationen, das Management von Drittkomponenten und Abhaengigkeiten, die Rolle von SBOMs fuer Komponenten-Transparenz und Patchbewertung, sowie Prozesse fuer Schwachstellenbehandlung, sichere Updates und Supportzeitraeume. Hinzu kommen technische Dokumentation, Risikobewertung und die Frage, welche Nachweise Hersteller vor dem Inverkehrbringen fuer Konformitaetsbewertung und CE-Kennzeichnung vorbereiten muessen. Ebenso relevant ist, wie Hersteller Incident Reporting und Meldungen zu aktiv ausgenutzten Schwachstellen organisatorisch und technisch vorbereiten, ohne den Embedded-Betrieb unnoetig zu verkomplizieren. Das Thema eignet sich besonders fuer eine Analyse typischer Zielkonflikte: begrenzte Ressourcen auf Mikrocontrollern, lange Einsatzdauer im Feld, OTA-Updatefaehigkeit, Integritaet der Lieferkette und Nachweisfuehrung gegenueber Marktaufsicht. Dadurch laesst sich der CRA als technischer Ordnungsrahmen fuer sichere vernetzte Produkte analysieren, statt nur als abstrakte Regulierung oder reine Compliance-Checkliste.}

\TopicLeitfragen{%
\item Welche CRA-Pflichten beeinflussen Architektur, Default-Konfiguration und Updatepfad vernetzter Embedded-Geraete unmittelbar?
\item Wie werden Supportzeitraum, Schwachstellenmanagement und sichere OTA-Updates bei ressourcenbeschraenkten IoT-Produkten technisch umgesetzt?
\item Welche Funktion erfuellen SBOMs fuer Komponenten-Transparenz, Patchbewertung und technische Dokumentation?
\item Wie veraendern Meldefristen fuer Vorfaelle und aktiv ausgenutzte Schwachstellen die Herstellerprozesse?
}

\TopicScope{%
\item Keine vertiefte juristische Auslegung von Bussgeldern, Haftungsfragen oder Vertragsgestaltung.
\item Keine vollstaendige Analyse aller harmonisierten Standards oder Common Specifications.
\item Kein Fokus auf reine Enterprise-Software ohne Embedded- oder IoT-Bezug.
\item Keine Schritt-fuer-Schritt-Compliance-Checkliste fuer ein konkretes Produkt.
}

\TopicPractice{Anhand eines typischen IoT-Geraets, z.~B. eines vernetzten Sensors oder Gateways, kann eine CRA-nahe Artefaktkette skizziert werden: Risikobewertung, Updatepfad, SBOM, Supportzeitraum und Meldeprozess. Dazu genuegt ein strukturiertes Architekturbeispiel ohne realen Hardware-Aufbau.}

\TopicKeywords{Cyber Resilience Act, Embedded Security, IoT-Sicherheit, Security by Design, Schwachstellenmanagement, OTA-Updates, SBOM, Incident Reporting}

\nocite{cra_embedded_iot_eurlex2024,cra_embedded_iot_ecmanufacturers2026,cra_embedded_iot_ecsummary2025,cra_embedded_iot_ecreport2026,cra_embedded_iot_jrc2024,cra_embedded_iot_enisa2026}
\end{topic}
